\section{Distributed Debugging Under EDT}
\label{sec:edt-debug-case}
This section gives a concrete, minimal workflow showing how EDT enables
\emph{distributed debugging} under a deterministic logical-time schedule. The key
idea is to pause at a safe window boundary, inspect which nodes are scheduled
next, and attach a native debugger to a specific client process (PID) while
preserving reproducibility.

\subsection{Start and pause at \(t=0\)}
When running interactively, EDT prints a small command set and auto-pauses at
\(t=0\) to provide a clean initial cut:

\begin{verbatim}
** Commands: c | cN (e.g. c10) | n | p | s | s:<pid> | info | r | rN
** [auto-paused] t=0s
\end{verbatim}

\subsection{Inspect the next window and select a target process}
\noindent\textbf{List next-window candidates.}
At a pause point, the \texttt{s} command enumerates hosts scheduled in the next
window, including each host's next-event time and the native PID(s) of its real
client process(es). This provides an immediate mapping from \emph{logical time}
to \emph{which real processes will execute next}.

\begin{verbatim}
s
[next-window] t=...s window=[...--...]
  host=geth-node          next_event=...   pid=3248123 (geth.1000)
  host=prysm-beacon-1     next_event=...   pid=3248456 (beacon-chain.1000)
  host=prysm-validator-1  next_event=...   pid=3248567 (validator.1000)
\end{verbatim}

\noindent\textbf{Optional: show more context.}
The \texttt{info} command prints the boundary context (e.g., current time,
window range, next global event), which is useful when aligning debugging with a
suspected failing interval:

\begin{verbatim}
info
[... paste your real terminal output here ...]
\end{verbatim}

\subsection{Attach a native debugger to one node}
\noindent\textbf{Get an attach command.}
Suppose we want to debug the geth node. We use \texttt{s:<pid>} to print a
copy-pastable attach command:

\begin{verbatim}
s:3248123
gdb -p 3248123
\end{verbatim}

\noindent\textbf{Why window boundaries enable ``distributed'' debugging.}
EDT pauses at \emph{window boundaries} that correspond to the runahead-style safe
horizon: within one window, hosts are guaranteed not to causally affect each
other via cross-host deliveries. Therefore, the tester can focus on one target
process (attach debugger, set breakpoints, single-step) while preserving a
consistent global cut and a deterministic schedule for the rest of the system.

\subsection{Step and continue deterministically}
\noindent\textbf{Step one window.}
The \texttt{n} command advances the system by exactly one window and pauses again
at the next boundary. This is the natural ``next'' operation because the window
is the minimal unit of safe parallel runahead:

\begin{verbatim}
n
[... paste your real output here ...]
\end{verbatim}

\noindent\textbf{Continue until a target time.}
To quickly reach a later region while keeping determinism, \texttt{cN} continues
until logical time \(t=N\) seconds and then pauses:

\begin{verbatim}
c120
[paused] t=120s
\end{verbatim}

\subsection{Lightweight replay via restart}
When a failure is observed (or after enabling extra instrumentation), \texttt{rN}
restarts from \(t=0\) and deterministically re-runs until \(t=N\) seconds,
pausing at the same boundary:

\begin{verbatim}
r120
** restarting simulation in-process
[paused] t=120s
\end{verbatim}

In practice, this workflow lets a tester reproduce and localize distributed
failures by iterating on a single logical-time window: pause, inspect, attach,
step, and replay, without relying on heavyweight snapshot mechanisms.
